% homework/hw07/hw07.tex
% Year 7 - Homework Packet 7: Unit 2, both subjects, and the place they meet.
%
% Covers Mathematics 2.3-2.4 (collecting like terms; expanding brackets) and
% Science 2.5-2.6 (atoms, elements and the Periodic Table; compounds and
% formulae).
%
% SECTION E IS THE POINT OF THIS PACKET, and the tie is exact rather than
% decorative. Both subjects handed this class the same two rules in the same
% week:
%
%   * a symbol says WHAT KIND, a number says HOW MANY;
%   * a number written in front multiplies EVERYTHING after it.
%
% From the first rule: 8s - s = 7s because s means 1s, and H2O has one oxygen
% atom because nothing is written after the O. Two of the same kind is still
% one kind, so x + x = 2x and O2 is still an element -- E4 is those two facts
% shown to be one fact.
%
% From the second: 3CO2 is three lots of CO2, which is 3(C + 2O) = 3C + 6O.
% That is 2.4's grid drawn on a chemical formula, and it is not an analogy --
% it is how the atoms in a formula are actually counted. E5 closes on the
% shared error: 4(x + 3) = 4x + 3 and 6CO2 = 6C + 2O are one mistake made
% twice, once in each subject.
%
% Nothing here is lifted from the workbook. Where a trap from a deck is reused
% it has new numbers and a new context.
%
% House style is shared: see homework/hw-style.tex. Method: the new-homework
% skill. Source is pure ASCII on purpose.
%
% Build:  pwsh .claude/skills/new-homework/scripts/build-hw.ps1 hw07

\documentclass[12pt,a4paper]{article}

% -- Answer key switch -------------------------------------------------------
\newif\ifkey
\ifdefined\TEACHER\keytrue\else
  \keyfalse        % \keyfalse = student copy   |   \keytrue = teacher copy
\fi

% -- House style -------------------------------------------------------------
% homework/hw-style.tex
% Shared house style for every Year 7 homework packet. \input this from the
% packet file, right after \documentclass and the answer-key switch.
%
% Do not fork this file per packet. If a packet needs a different look, the
% question is whether every packet should have it.
%
% The choices here are deliberate and were arrived at by printing and looking:
%   * No solid dark banners. Every header is a light tint with a coloured spine
%     on the left, so colour reads clearly without flooding a school printer.
%   * No marks and no mark totals. These are practice, not exam papers.
%   * Lato at 12pt with open leading, plus mathastext so inline maths is Lato
%     too. Without mathastext every "-6 + 4" is Computer Modern serif and the
%     sheet looks half-typeset.
%   * Writing lines are spaced for Year 7 handwriting, not adult margin notes.
%
% Provides: \wlines \blank \tickbox \sep \qmk-free marks (none) \hwsection
% \qhead \expr, the workspace box, and the remember / wordhelp / activity /
% yourturn coloured boxes. Palette matches the lesson decks exactly.

% -- Packages ----------------------------------------------------------------
\usepackage[T1]{fontenc}
\usepackage[utf8]{inputenc}
\usepackage[default]{lato}
\usepackage[a4paper,top=1.7cm,bottom=1.7cm,left=2.0cm,right=2.0cm,headheight=15pt]{geometry}
\usepackage{amsmath,amssymb}
\usepackage{xcolor}
\usepackage[most]{tcolorbox}
\usepackage{enumitem}
\usepackage{tikz}
\usetikzlibrary{arrows.meta}
\usepackage{array}
\usepackage{tabularx}
\usepackage{fancyhdr}
\usepackage{multicol}
\usepackage{needspace}
\usepackage{graphicx}

% Make maths use Lato too. Without this, every "-6 + 4" on the page is set in
% Computer Modern serif and the sheet looks half-typeset. Must come last.
\usepackage[italic]{mathastext}

\linespread{1.06}\selectfont

% -- The Learner's Book palette (same hex values as the lesson decks) ---------
\definecolor{cbteal}{HTML}{0087A8}
\definecolor{cbpurple}{HTML}{5C2483}
\definecolor{cborange}{HTML}{C25E12}
\definecolor{cbgreen}{HTML}{4A8B23}
\definecolor{cbred}{HTML}{C8102E}
\definecolor{cbblue}{HTML}{1A5FA8}
\definecolor{rulegray}{HTML}{9AA5AD}

% -- Answer lines and blanks -------------------------------------------------
% \wlines{n} draws n ruled writing lines, spaced wide enough for Year 7
% handwriting rather than for an adult with a fine pen.
\makeatletter
\newcommand{\wlines}[1]{%
  \par\vspace{0.75em}%
  \@tempcnta=0\relax
  \@whilenum\@tempcnta<#1\do{%
    \hbox to \linewidth{\color{rulegray}\leaders\hrule height 0.5pt\hfill}%
    \vspace{1.5em}%
    \advance\@tempcnta by 1\relax}%
  \par\vspace{-0.8em}%
}
\makeatother

% \blank[width] is an inline gap to write a short answer in.
\newcommand{\blank}[1][2.6cm]{\underline{\hspace{#1}}}

% A boxed space for showing working.
\newtcolorbox{workspace}[1][3.4cm]{%
  colback=white, colframe=rulegray, boxrule=0.5pt, arc=5pt,
  height=#1, left=6pt, right=6pt, top=4pt, bottom=4pt,
  before skip=9pt, after skip=11pt}

% An empty tick box.
\newcommand{\tickbox}{\fbox{\rule{0pt}{1.15em}\rule{1.15em}{0pt}}}

% A middot separator.
\newcommand{\sep}{\,\textperiodcentered\,}

% -- Section and question headings -------------------------------------------
% Light tint plus a fat coloured spine on the left. No solid dark banner: the
% colour reads clearly but barely costs the printer anything.
% needspace stops a heading stranding alone at the foot of a page. Kept modest:
% too large and it pushes headings over, leaving a worse hole than it prevents.
\newcommand{\hwsection}[2][cbteal]{%
  \needspace{8\baselineskip}%
  \par\vspace{13pt}%
  \noindent
  \begin{tcolorbox}[enhanced, nobeforeafter, width=\textwidth,
      colback=#1!7, colframe=#1!7, arc=5pt,
      borderline west={5pt}{0pt}{#1},
      left=14pt, right=10pt, top=7pt, bottom=7pt]
    {\large\bfseries\color{#1}#2}
  \end{tcolorbox}%
  \par\vspace{9pt}}

% \qhead[colour]{A2}{Moving along the line} -- a question heading.
\newcommand{\qhead}[3][cbteal]{%
  \needspace{6\baselineskip}%
  \par\vspace{11pt}%
  \noindent{\large\bfseries\color{#1}#2}\quad{\large\bfseries #3}\par\vspace{4pt}}

% -- Coloured boxes ----------------------------------------------------------
% Shared look: rounded, light fill, coloured rule, coloured title text on a
% pale title strip rather than a dark one.
\tcbset{hwbox/.style={
  enhanced, breakable, arc=6pt, boxrule=1pt,
  left=11pt, right=11pt, top=8pt, bottom=8pt,
  fonttitle=\bfseries, attach boxed title to top left={xshift=12pt, yshift=-8pt},
  boxed title style={arc=4pt, boxrule=1pt, left=7pt, right=7pt, top=2pt, bottom=2pt},
  top=13pt}}

% Orange: things to remember. Matches the deck's "Write This Down".
\newtcolorbox{remember}[1][Remember from the lesson]{hwbox,
  colback=cborange!6, colframe=cborange, coltitle=cborange,
  colbacktitle=white, title={#1}}

% Blue: English help. The barrier in this class is the wording, not the maths.
\newtcolorbox{wordhelp}[1][Word help]{hwbox,
  colback=cbblue!6, colframe=cbblue, coltitle=cbblue,
  colbacktitle=white, title={#1}}

% Purple: an activity or instruction.
\newtcolorbox{activity}[1][How to do this homework]{hwbox,
  colback=cbpurple!6, colframe=cbpurple, coltitle=cbpurple,
  colbacktitle=white, title={#1}}

% Green: your turn / a friendly nudge.
\newtcolorbox{yourturn}[1][Your turn]{hwbox,
  colback=cbgreen!7, colframe=cbgreen, coltitle=cbgreen!75!black,
  colbacktitle=white, title={#1}}

% -- Shorthands --------------------------------------------------------------
\newcommand{\dC}{\,^{\circ}\mathrm{C}}          % use inside maths: $-8\dC$
\newcommand{\key}[1]{{\color{cborange}\textbf{#1}}}
\newcommand{\eng}[1]{{\color{cbblue}\textbf{#1}}}

\setlist[enumerate]{itemsep=8pt, topsep=5pt, leftmargin=*}
\setlist[itemize]{itemsep=4pt, topsep=4pt, leftmargin=*}
\renewcommand{\arraystretch}{1.45}

% -- An expression the student is asked to work from -------------------------
\newcommand{\expr}[1]{%
  \tcbox[on line, colback=cbgreen!10, colframe=cbgreen!55, boxrule=1pt, arc=4pt,
         left=9pt, right=9pt, top=4pt, bottom=4pt]{\large\bfseries $#1$}}

% -- Running head and foot ---------------------------------------------------
% The packet overrides \fancyhead[L] with its own title.
\pagestyle{fancy}
\fancyhf{}
\fancyhead[L]{\small\color{black!55}Year 7\sep Homework}
\fancyhead[R]{\small\color{black!55}Name: \blank[3.4cm]}
\fancyfoot[C]{\small\color{black!55}Page \thepage}
\renewcommand{\headrulewidth}{0pt}
\renewcommand{\headrule}{\vspace{2pt}\hbox to\headwidth{\color{cbteal!45}\leaders\hrule height 1.2pt\hfill}}



% -- Per-packet running head -------------------------------------------------
\fancyhead[L]{\small\color{black!55}Year 7\sep Homework 7}

% -- A group of writing lines that will not be split -------------------------
% See hw06.tex: \wlines is a run of separate hboxes and TeX will break in the
% middle of one, stranding a single rule at the top of the next page.
\newcommand{\qlines}[1]{\par\needspace{\numexpr#1+1\relax\baselineskip}\wlines{#1}}

% -- A chemical formula ------------------------------------------------------
% Upright, so it does not come out in mathastext's italic and read as algebra.
% That distinction is the whole of E1, so it has to be visibly true everywhere
% else on the packet too.
\newcommand{\ch}[1]{\ensuremath{\mathrm{#1}}}

% -- A box to write one number in --------------------------------------------
\newcommand{\nbox}{\ \fbox{\rule{0pt}{1.05em}\rule{1.0em}{0pt}}\ }

% -- An answer blank that will not cause an overfull line ---------------------
% "\hfill \blank[3.4cm]" at the end of a wrapping \item overflows the margin
% when the text happens to fill its last line: the blank has nowhere to go and
% hangs off the right edge. Two of them did. \ansline puts the blank on a line
% of its own, right-aligned, which cannot overflow whatever the text does.
\newcommand{\ansline}[1][4.4cm]{\par\vspace{2pt}\hspace*{\fill}\blank[#1]\par}

% -- A note on the word-help tables ------------------------------------------
% They are written {@{}l@{\hspace{16pt}}X@{}}, spelled out in full. The first
% column must be `l', not `p': LaTeX puts the stretched row strut in the FIRST
% cell only, and inside a p-column it lands in the parbox and pushes the text
% down a line, so every left-hand phrase prints below its meaning. The spec
% cannot be hidden behind a macro (array does not expand macros in a preamble)
% nor the table behind a wrapper environment (tabularx reads ahead for a
% literal end-of-tabularx). Both were tried.

\begin{document}

% ===========================================================================
% COVER BLOCK
% ===========================================================================
\thispagestyle{fancy}

\begin{tcolorbox}[enhanced, colback=cbpurple!7, colframe=cbpurple!7, arc=7pt,
                  borderline west={6pt}{0pt}{cbpurple},
                  left=18pt, right=14pt, top=10pt, bottom=10pt]
  {\color{cbpurple}\bfseries YEAR 7\sep UNIT 2\sep HOMEWORK 7}\\[5pt]
  {\color{cbpurple}\Huge\bfseries What Kind, and How Many}\\[7pt]
  {\color{black!70}Mathematics 2.3--2.4 --- collecting like terms, expanding brackets}\\
  {\color{black!70}Science 2.5--2.6 --- atoms, elements, compounds and formulae}\\
  {\color{black!70}Section E --- the two rules both subjects taught you this week}
\end{tcolorbox}

\vspace{4pt}
\noindent
\begin{tabularx}{\textwidth}{@{}X X@{}}
  \textbf{Name:} \blank[5.6cm] & \textbf{Class:} \blank[5.6cm] \\
  \textbf{Date given:} \blank[4.6cm] &
  \textbf{Date due:} {\color{cbred}\bfseries Monday 21 September 2026} \\
\end{tabularx}

\vspace{2pt}
\begin{activity}
  Five sections. Maths first, then science, then \key{Section E}, where the two
  meet.
  \begin{itemize}
    \item In maths, $3x + 5$ is a \key{finished answer}. You are allowed to stop.
    \item In science, write \key{full English sentences}. ``Compound'' is not a
          sentence; ``Salt is a compound of sodium and chlorine'' is.
    \item \eng{No calculator.}
  \end{itemize}
\end{activity}

% ===========================================================================
% SECTION A -- MATHEMATICS 2.3
% ===========================================================================
\hwsection{Section A \textperiodcentered\ Mathematics 2.3 --- Collecting Like Terms}

% -- A1 ---------------------------------------------------------------------
% The way in, and it is meant to be easy: count the terms. A student who can
% start nothing else can start this. The instruction line comes BEFORE the
% orange box -- a section heading followed straight by a large tcolorbox cannot
% be split, and the pair walks to the next page and leaves a hole behind it.
\qhead{A1}{How many terms?}

\noindent Count the terms in each expression. The first one is done for you.

\begin{remember}[What a term is]
  A \key{term} is one part of an expression. The \key{plus and minus signs
  separate the terms}. A number on its own is a term too. \; So $4a + 3b$ has
  \key{two} terms: $4a$ and $3b$.
\end{remember}

\vspace{4pt}
\noindent
\begin{tabularx}{\textwidth}{|p{3.9cm}|X|p{3.9cm}|X|}
  \hline
  \textbf{Expression} & \textbf{Terms} & \textbf{Expression} & \textbf{Terms} \\
  \hline
  $4a + 3b$ \rule[-0.65em]{0pt}{2.0em}      & \textbf{2}
    & $8m - 3m + n$ \rule[-0.65em]{0pt}{2.0em} & \\ \hline
  $5x$ \rule[-0.65em]{0pt}{2.0em}           &
    & $6$ \rule[-0.65em]{0pt}{2.0em}           & \\ \hline
  $2p + 7q + 4$ \rule[-0.65em]{0pt}{2.0em}  &
    & $y + 2z - 5 + 3y$ \rule[-0.65em]{0pt}{2.0em} & \\ \hline
\end{tabularx}

% -- A2 ---------------------------------------------------------------------
% One line of prose before the word-help box, for the same reason as A1.
\qhead{A2}{Like or not like?}

\noindent Tick one box in each row.

\begin{wordhelp}[Word help --- one word with two jobs]
  \begin{tabularx}{\linewidth}{@{}l@{\hspace{16pt}}X@{}}
    \eng{like} & every day, to enjoy something. In maths, \textbf{the same kind}: \emph{milk is like water --- both are liquids.} \\
    \eng{like terms} & terms with \textbf{the same letter}, so they can be added together \\
  \end{tabularx}
\end{wordhelp}

\vspace{2pt}
\noindent
\begin{tabularx}{\textwidth}{|X|c|c|X|c|c|}
  \hline
  & \textbf{Like} & \textbf{Not like} & & \textbf{Like} & \textbf{Not like} \\
  \hline
  $3a$ and $5a$ \rule[-0.65em]{0pt}{2.0em}   & & &
  $2ab$ and $5ba$ \rule[-0.65em]{0pt}{2.0em} & & \\ \hline
  $4x$ and $4y$ \rule[-0.65em]{0pt}{2.0em}   & & &
  $m$ and $m^2$ \rule[-0.65em]{0pt}{2.0em}   & & \\ \hline
  $7$ and $2$ \rule[-0.65em]{0pt}{2.0em}     & & &
  $5c$ and $-2c$ \rule[-0.65em]{0pt}{2.0em}  & & \\ \hline
\end{tabularx}

% -- A3 ---------------------------------------------------------------------
% The bare-calculation block, with the invisible 1 inside it rather than in a
% question of its own. Items (e)-(h) each contain a lone letter, so the rule in
% the box is needed four times and never announced again.
\qhead{A3}{Simplify}

\noindent Collect the like terms. Two of these are already finished --- leave them alone.

\begin{remember}[$x$ means $1x$]
  A letter on its own already has a \key{1} in front of it. Nobody writes the 1,
  but it is there.

  \smallskip
  \hspace*{1.2em}$8s - s \;=\; 8s - 1s \;=\; \mathbf{7s}$ \qquad
  {\color{cbred}\bfseries not 8.}
\end{remember}

\vspace{2pt}
\begin{multicols}{2}
\begin{enumerate}[label=(\alph*), itemsep=10pt]
  \item $2b + 3b = \blank[2.4cm]$
  \item $9k - 4k = \blank[2.4cm]$
  \item $6m + 3n + 2m = \blank[2.4cm]$
  \item $5p + 2q + p + 4q = \blank[2.4cm]$
  \item $8s - s = \blank[2.4cm]$
  \item $4x + x + 2x = \blank[2.4cm]$
  \item $3a + 7 = \blank[2.4cm]$
  \item $6y - y + 2y = \blank[2.4cm]$
  \item $10g - 6g - g = \blank[2.4cm]$
  \item $5c + 3d = \blank[2.4cm]$
\end{enumerate}
\end{multicols}

% -- A4 ---------------------------------------------------------------------
\qhead{A4}{Every sign travels with its own term}

\begin{remember}[Take the sign with you]
  The $+$ or $-$ in front of a term \key{belongs to that term}. When you move the
  term, the sign goes with it.

  \smallskip
  \hspace*{1.2em}$7x + 5y - 3x + y \;=\; 7x - 3x + 5y + y \;=\; \mathbf{4x + 6y}$
\end{remember}

\noindent Simplify. Write the middle line, with the like terms next to each other.

\vspace{4pt}
\begin{enumerate}[label=(\alph*), itemsep=10pt]
  \item $9x + 4y - 5x + 2y$ \hfill \blank[6.6cm]
  \item $8a - 3b + 2a - b$ \hfill \blank[6.6cm]
  \item $6 + 5t - 1 - 2t$ \hfill \blank[6.6cm]
  \item $7p - 2q - 3p + 6q + 4$ \hfill \blank[6.6cm]
\end{enumerate}

% -- A5 ---------------------------------------------------------------------
% There was a "Mr Bowen's homework" error-hunt here -- 4x + 3 = 7x, 7y - y = 7,
% 3p + 2q + p, 4ab + 2ba. It came out for the page target and it was the right
% one to cut, because every one of its four traps is already on the packet:
% 4ab/2ba is a row of A2, the two already-finished items are A3(g) and A3(j),
% and the invisible 1 is four separate items of A3. B4 keeps the error-hunt
% format, on 2.4, where it pairs with E5.
%
% The reverse question: an expression, and the student writes the word problem.
% It cannot be pattern-matched, so it is the best single item in Section A for
% finding out who is reading.
\qhead{A5}{Now you write the question}

\begin{yourturn}[Read all four conditions before you start]
  Here is an expression:\qquad \expr{5a + 3b}

  \medskip
  Write a \key{word problem} that this expression is the answer to. Your question
  must:
  \begin{enumerate}[label=\arabic*., itemsep=4pt, topsep=5pt]
    \item be about something \key{real} --- a shop, a classroom, a lab bench;
    \item use \key{two different kinds} of thing, one for $a$ and one for $b$;
    \item say clearly what $a$ and $b$ \key{stand for};
    \item be written in \key{full sentences} and \key{end in a question}.
  \end{enumerate}
\end{yourturn}

\qlines{4}

% ===========================================================================
% SECTION B -- MATHEMATICS 2.4
% ===========================================================================
\hwsection{Section B \textperiodcentered\ Mathematics 2.4 --- Expanding Brackets}

% -- B1 ---------------------------------------------------------------------
\qhead{B1}{Fill the grid}

\noindent One box, one multiplication. Fill in the empty boxes, then write the answer.

\begin{wordhelp}[Word help --- the words that carry the whole rule]
  \begin{tabularx}{\linewidth}{@{}l@{\hspace{16pt}}X@{}}
    \eng{expand} \sep \eng{multiply out} & the same instruction in different words. Both appear in the exercise \\
    \eng{each / every} & not a maths word, but it carries the whole rule: \textbf{every} term inside, not just the first \\
  \end{tabularx}
\end{wordhelp}

\begin{remember}[Brackets mean multiply]
  $4(x + 3)$ means $4 \times (x + 3)$. The times sign is hidden, exactly as it is
  in $4n$.

  \medskip
  \noindent
  \begin{minipage}[c]{0.40\linewidth}
    \renewcommand{\arraystretch}{1.6}%
    \begin{tabular}{|c|c|c|}
      \hline
      $\times$ & $x$  & $+\,3$ \\ \hline
      $4$      & $4x$ & $+\,12$ \\ \hline
    \end{tabular}
  \end{minipage}%
  \begin{minipage}[c]{0.58\linewidth}
    $4(x + 3) = \mathbf{4x + 12}$

    \smallskip
    \eng{Not} $4x + 3$. The 3 is inside the brackets, so it gets multiplied too.
  \end{minipage}
\end{remember}

% Three separate blocks with a real \vspace between them, NOT three rows of one
% tabularx. As rows, the grids stack with their borders touching and read as a
% single six-row table -- which is exactly the wrong idea, since each grid is its
% own multiplication. The optional argument of \\ does not open the gap here:
% each inner tabular is centred on its row's baseline and swallows it. Explicit
% blocks are the only thing that reliably separates them.
\renewcommand{\arraystretch}{1.55}%
\vspace{4pt}

\noindent
\begin{minipage}[c]{0.26\textwidth}
  \begin{tabular}{|c|c|c|}
    \hline
    $\times$ & $m$ & $+\,4$ \\ \hline
    $3$      &     &        \\ \hline
  \end{tabular}
\end{minipage}%
\begin{minipage}[c]{0.72\textwidth}
  \textbf{(a)} \; $3(m + 4) = \blank[4.0cm]$
\end{minipage}

\vspace{13pt}
\noindent
\begin{minipage}[c]{0.26\textwidth}
  \begin{tabular}{|c|c|c|}
    \hline
    $\times$ & $y$ & $+\,2$ \\ \hline
    $5$      &     &        \\ \hline
  \end{tabular}
\end{minipage}%
\begin{minipage}[c]{0.72\textwidth}
  \textbf{(b)} \; $5(y + 2) = \blank[4.0cm]$
\end{minipage}

\vspace{13pt}
\noindent
\begin{minipage}[c]{0.26\textwidth}
  \begin{tabular}{|c|c|c|}
    \hline
    $\times$ & $2n$ & $+\,3$ \\ \hline
    $4$      &      &        \\ \hline
  \end{tabular}
\end{minipage}%
\begin{minipage}[c]{0.72\textwidth}
  \textbf{(c)} \; $4(2n + 3) = \blank[4.0cm]$
\end{minipage}
\renewcommand{\arraystretch}{1.45}%

% -- B2 ---------------------------------------------------------------------
% Plain expanding, a minus inside, and a number in front of the letter, all in
% one block. Three headings and three boxes for what is really one skill left
% the section a page longer than the packet has.
\qhead{B2}{Expand}

\noindent No grid this time. Draw one if it helps.

\begin{remember}[Two things that go inside with the term]
  \key{The minus sign.} \; $3(x - 2) = 3x - 6$ \quad and \quad $7(1 - w) = 7 - 7w$.

  \smallskip
  \key{A number already in front of the letter.} Multiply the numbers, keep the
  letter: \; $5 \times 2p = 10p$, \; {\color{cbred}\bfseries not $52p$}.
\end{remember}

\vspace{2pt}
\begin{multicols}{3}
\begin{enumerate}[label=(\alph*), itemsep=11pt]
  \item $3(a + 2) = \blank[1.5cm]$
  \item $5(b + 3) = \blank[1.5cm]$
  \item $9(y + 3) = \blank[1.5cm]$
  \item $4(2 + f) = \blank[1.5cm]$
  \item $8(7 + z) = \blank[1.5cm]$
  \item $6(k - 3) = \blank[1.5cm]$
  \item $2(y - 4) = \blank[1.5cm]$
  \item $5(m - 6) = \blank[1.5cm]$
  \item $4(3 - t) = \blank[1.5cm]$
  \item $3(2q + 5) = \blank[1.5cm]$
  \item $4(5u - 1) = \blank[1.5cm]$
  \item $6(1 + 2v) = \blank[1.5cm]$
\end{enumerate}
\end{multicols}

\vspace{2pt}
\noindent \textbf{(m)} This one has three terms inside. \quad $8(6 + 4w - 3g) =
\blank[5.0cm]$

% -- B3 ---------------------------------------------------------------------
% 2.3 arriving inside 2.4. This is why Section A comes first.
\qhead{B3}{Knowing when to stop}

\noindent Tick one box in each row. If it can be collected, write the answer.

\begin{remember}[Expanding does not switch off Section A]
  After you expand, look for like terms. Often there are none.

  \smallskip
  \hspace*{1.2em}$4(3 - c) \;=\; 12 - 4c$, \; and $12$ and $4c$ are
  \key{not like terms}. So $12 - 4c$ is \key{finished}.
\end{remember}

\vspace{2pt}
\noindent
\begin{tabularx}{\textwidth}{|p{3.7cm}|c|c|X|}
  \hline
  \textbf{Expression} & \textbf{Finished} & \textbf{Can be collected} &
  \textbf{If so, write it} \\
  \hline
  $12 - 4c$ \rule[-0.7em]{0pt}{2.2em}       & & & \\ \hline
  $5x + 3x$ \rule[-0.7em]{0pt}{2.2em}       & & & \\ \hline
  $3a + 6$ \rule[-0.7em]{0pt}{2.2em}        & & & \\ \hline
  $8 + 2m + 5$ \rule[-0.7em]{0pt}{2.2em}    & & & \\ \hline
\end{tabularx}

% -- B4 ---------------------------------------------------------------------
\qhead{B4}{Mr Bowen's homework, again}

\noindent One of these four is right. Find it, and correct the other three.

\vspace{4pt}
\begin{multicols}{2}
\begin{enumerate}[label=(\alph*), itemsep=10pt]
  \item $3(a + 4) = 3a + 4$ \; \blank[2.2cm]
  \item $4(2b - 5) = 8b - 9$ \; \blank[2.2cm]
  \item $2(5 - d) = 10 - 2d$ \; \blank[2.2cm]
  \item $6(m + 2) = 6m + 8$ \; \blank[2.2cm]
\end{enumerate}
\end{multicols}

% -- B5 ---------------------------------------------------------------------
\qhead{B5}{Expand, then collect}

\noindent Two jobs, in this order. Expand the brackets first, then look for like terms.

\vspace{2pt}
\begin{enumerate}[label=(\alph*), itemsep=4pt]
  \item $3(2y + 1) + 5y + 4$
        \begin{workspace}[1.7cm]\end{workspace}
  \item $2(a + 6) + 3(a + 1)$
        \begin{workspace}[1.7cm]\end{workspace}
\end{enumerate}

% -- B6 ---------------------------------------------------------------------
% There was a "work backwards" question here -- fill the box in [](x + 3) =
% 5x + 15 -- which is the deck's Challenge slide. It came out for the page
% target and it is the right thing to lose, because it is the only item in
% Section B that is not one of 2.4's stated objectives: B1 to B5 cover expand,
% the sign, the number in front, when to stop, and expand-and-simplify between
% them. The skill itself is not lost to the class -- hw07x teaches it at length,
% as the place where dividing turns out to be what undoes a bracket.
\qhead{B6}{The fish tanks}

\noindent Mr Bowen keeps \textbf{7} fish tanks in his kitchen. Every tank holds $f$
fish and \textbf{2} rubber ducks.

\vspace{4pt}
\begin{enumerate}[label=(\alph*), itemsep=9pt]
  \item Write an expression with brackets for the number of things in the tanks
        altogether. \ansline
  \item Expand it. \hfill \blank[4.4cm]
  \item Rubber ducks are not alive. Write an expression for the number of
        \textbf{living} things in the tanks.
        \ansline
\end{enumerate}

% ===========================================================================
% SECTION C -- SCIENCE 2.5
% ===========================================================================
\hwsection{Section C \textperiodcentered\ Science 2.5 --- Atoms, Elements and the Periodic Table}

% -- C1 ---------------------------------------------------------------------
% This question opens the section because it is small, and because the section
% header lands near the foot of a page: C2's eight-row table needs about 10 cm
% and will not follow it, so without something short here the page ends a fifth
% white. A two-row lookup fits in the gap exactly.
%
% It is also the only question on the packet that cannot be answered without
% opening the book, which is half of what 2.5 was about. A five-element version
% was cut for the page target; this is what came back.
\qhead{C1}{Use the Periodic Table in your book}

\noindent Groups are numbered 1 to 8.

\vspace{6pt}
\noindent
\begin{tabularx}{\textwidth}{|X|p{2.0cm}|p{2.0cm}|p{2.0cm}|p{3.1cm}|}
  \hline
  \textbf{Element} & \textbf{Symbol} & \textbf{Period} & \textbf{Group} &
  \textbf{Metal or} \newline \textbf{non-metal} \\
  \hline
  magnesium \rule[-0.7em]{0pt}{2.2em} & & & & \\ \hline
  chlorine \rule[-0.7em]{0pt}{2.2em}  & & & & \\ \hline
\end{tabularx}

% -- C2 ---------------------------------------------------------------------
\qhead{C2}{Match the word to its meaning}

\noindent Write the correct \textbf{letter} in each box.

\vspace{7pt}
\noindent
\begin{tabularx}{\textwidth}{|p{3.9cm}|c|X|}
  \hline
  \textbf{Word} & & \textbf{Meaning} \\
  \hline
  1.\ Atom              & \tickbox & \textbf{A} \; A column in the Periodic Table. \\
  2.\ Element           & \tickbox & \textbf{B} \; A short way to write the name of an element. \\
  3.\ The Periodic Table & \tickbox & \textbf{C} \; Elements such as iron, copper and aluminium. \\
  4.\ Period            & \tickbox & \textbf{D} \; A tiny piece of matter. Everything is made of them. \\
  5.\ Group             & \tickbox & \textbf{E} \; A very, very small tube made of carbon atoms. \\
  6.\ Symbol            & \tickbox & \textbf{F} \; A substance made of only one kind of atom. \\
  7.\ Nanotube          & \tickbox & \textbf{G} \; A row in the Periodic Table. \\
  8.\ Metals            & \tickbox & \textbf{H} \; A way of arranging all the elements. \\
  \hline
\end{tabularx}

% -- C3 ---------------------------------------------------------------------
\qhead{C3}{The capital-letter rule}

\noindent \textbf{(a)} Four students wrote a symbol wrongly. Write it correctly, and
name the element.

\begin{remember}[One capital, then small letters]
  An element symbol starts with a \key{CAPITAL} letter. If there is a second
  letter it is \key{small}. \; \ch{Mg} is magnesium --- \key{one} element.
  \ch{MG} and \ch{mg} are not symbols at all.
\end{remember}

\vspace{2pt}
\noindent
% Two self-contained three-column blocks side by side, not one row carrying two
% symbols. The earlier version put "NA / cl" in a single cell followed by four
% answer cells, and nothing said which Correct column belonged to which symbol.
\begin{tabularx}{\textwidth}{|p{1.9cm}|p{1.9cm}|X|p{1.9cm}|p{1.9cm}|X|}
  \hline
  \textbf{Written} & \textbf{Correct} & \textbf{Element} &
  \textbf{Written} & \textbf{Correct} & \textbf{Element} \\
  \hline
  \ch{NA} \rule[-0.7em]{0pt}{2.2em} & & & \ch{mG} \rule[-0.7em]{0pt}{2.2em} & & \\ \hline
  \ch{cl} \rule[-0.7em]{0pt}{2.2em} & & & \ch{CA} \rule[-0.7em]{0pt}{2.2em} & & \\ \hline
\end{tabularx}

\vspace{9pt}
\noindent \textbf{(b)} Mr Bowen means cobalt, and writes \ch{CO}. Cobalt is \ch{Co}.
Explain what he has actually written down.
\qlines{2}

% -- C4 ---------------------------------------------------------------------
% There was a part (b) here -- Mr Bowen halving a piece of gold 25 times down to
% one atom, and why he cannot halve it again. It went for the page target. C1's
% first row already gives the class what an atom is, and 2.5's stated objective
% is met by (a), which is the harder question of the two by some distance.
\qhead{C4}{Answer in full sentences}

\noindent Diamond is hard and clear. Graphite is soft and grey. Both are made of
\textbf{only} carbon atoms. Is diamond an element? Explain your answer.

\begin{wordhelp}[Sentence starters --- you may copy these]
  ``An element is a substance that\ldots'' \qquad
  ``\ldots because every atom in it is\ldots''
\end{wordhelp}

\qlines{3}

% ===========================================================================
% SECTION D -- SCIENCE 2.6
% ===========================================================================
\hwsection{Section D \textperiodcentered\ Science 2.6 --- Compounds and Formulae}

% -- D1 ---------------------------------------------------------------------
\qhead{D1}{Element or compound?}

\noindent Tick one box in each row.

\begin{remember}[Count the kinds, not the atoms]
  An \key{element} is made of \key{one kind} of atom. A \key{compound} is made of
  \key{two or more kinds} of atom \key{bonded} together. \; \ch{O_2} has two
  atoms and they are \key{both oxygen} --- one kind, so it is an \key{element}.
\end{remember}

\vspace{2pt}
\noindent
\begin{tabularx}{\textwidth}{|X|c|c|X|c|c|}
  \hline
  & \textbf{Element} & \textbf{Compound} & & \textbf{Element} & \textbf{Compound} \\
  \hline
  \ch{Na} \rule[-0.65em]{0pt}{2.0em}   & & & \ch{H_2} \rule[-0.65em]{0pt}{2.0em}    & & \\ \hline
  \ch{NaCl} \rule[-0.65em]{0pt}{2.0em} & & & \ch{CaCO_3} \rule[-0.65em]{0pt}{2.0em} & & \\ \hline
  \ch{O_2} \rule[-0.65em]{0pt}{2.0em}  & & & \ch{Al} \rule[-0.65em]{0pt}{2.0em}     & & \\ \hline
  \ch{CO_2} \rule[-0.65em]{0pt}{2.0em} & & & \ch{H_2O} \rule[-0.65em]{0pt}{2.0em}   & & \\ \hline
\end{tabularx}

% -- D2 ---------------------------------------------------------------------
% There was a "Read the formula" table here: H2O, CO2, CH4, NaOH, CaCO3, with
% columns for which elements and how many atoms in total. It came out for the
% page target. Its skill is not lost -- E3's middle column asks for the atoms of
% each element by name, which is the same reading done harder, and the one trap
% that lived only here, NaOH being sodium rather than nitrogen and hydrogen, is
% now a row of E1, where the capital-letter rule that decides it is the subject
% of the question.
\qhead{D2}{Naming a compound}

\noindent Name the compound made from each set of elements.

\begin{wordhelp}[Word help --- the endings and the counting words]
  \begin{tabularx}{\linewidth}{@{}l@{\hspace{16pt}}X@{}}
    \eng{-ide}  & two elements joined. The metal comes first; the non-metal ends in \textbf{-ide} \\
    \eng{-ate}  & two elements and \textbf{oxygen} joined. The ending is often \textbf{-ate} \\
    \eng{mono / di} & one / two. \ch{CO} is carbon mon\textbf{o}xide; \ch{CO_2} is carbon \textbf{di}oxide \\
  \end{tabularx}
\end{wordhelp}

\vspace{2pt}
\begin{multicols}{2}
\begin{enumerate}[label=(\alph*), itemsep=10pt]
  \item sodium and chlorine \\ \blank[4.6cm]
  \item magnesium and oxygen \\ \blank[4.6cm]
  \item lithium and sulfur \\ \blank[4.6cm]
  \item calcium, carbon and oxygen \\ \blank[4.6cm]
\end{enumerate}
\end{multicols}

\vspace{4pt}
\noindent \textbf{(e)} Mr Bowen writes \emph{sulfur calcide}. Two things are wrong.
Write the correct name and say what the two mistakes were.
\qlines{2}

% -- D4 ---------------------------------------------------------------------
\qhead{D3}{A compound is nothing like the elements in it}

\noindent Sodium is a soft metal that bursts into flame in water. Chlorine is a
poisonous yellow-green gas. Bonded together they make sodium chloride, which you ate
with your dinner last night.

\vspace{5pt}
\noindent A student says: \emph{``Salt must be dangerous, because chlorine is
poisonous.''} Explain, in full sentences, why this is wrong.
\qlines{3}

% -- D5 ---------------------------------------------------------------------
% Atomic block: four drawing panels and their labels cannot be split, so it sits
% at the end of Section D where it strands nothing.
\qhead{D4}{Draw the particles}

\noindent Draw the particles of each substance in the box. Use a different colour or
pattern for each kind of atom, then tick element or compound.

\vspace{7pt}
\noindent
\renewcommand{\arraystretch}{1.3}%
\begin{tabularx}{\textwidth}{|X|X|X|X|}
  \hline
  \textbf{oxygen} & \textbf{water} & \textbf{carbon dioxide} & \textbf{methane} \\
  \ch{O_2} & \ch{H_2O} & \ch{CO_2} & \ch{CH_4} \\
  \hline
  \rule[-5.2em]{0pt}{5.6em} & & & \\
  \hline
  element \tickbox & element \tickbox & element \tickbox & element \tickbox \\
  compound \tickbox & compound \tickbox & compound \tickbox & compound \tickbox \\
  \hline
\end{tabularx}
\renewcommand{\arraystretch}{1.45}%

% ===========================================================================
% SECTION E -- THE BRIDGE
% ===========================================================================
\hwsection[cbgreen]{Section E \textperiodcentered\ Where the two subjects meet}

\noindent In maths a letter stands for a number. In science a letter stands for an
element. You read them with the same two rules.

% The worked examples that used to sit under each rule came out. They were the
% examples E2 and E3 open with anyway, and printing them twice cost the bottom
% of this page: E1's heading, box and table are about 12 cm and would not fit
% under the longer version, so the page ended two-fifths white.
\begin{remember}[The two rules]
  \key{1. \; A letter says what kind. A number says how many.} \quad $3a$ \sep \ch{O_3}

  \smallskip
  \key{2. \; A number written in front multiplies everything after it.} \quad
  $4(x+3)$ \sep \ch{3CO_2}
\end{remember}

% -- E1 ---------------------------------------------------------------------
\qhead[cbgreen]{E1}{$\mathrm{Mg}$ or $mg$?}

\noindent Complete the table. The first row is done for you.

\begin{remember}[The capital is doing a job]
  In \key{maths} every letter is small, and two letters written together means
  \key{multiply}: \; $mg$ means $m \times g$.

  \smallskip
  In \key{science} a symbol begins with a \key{capital}: \; \ch{Mg} is magnesium,
  one element. Two capitals means \key{two} elements: \; \ch{CO} is carbon and
  oxygen.
\end{remember}

% The NaOH row arrived here when D2 was cut. It belongs here: what makes it hard
% is precisely the capital-letter rule, because Na is one symbol of two letters
% and every student who reads it as nitrogen-and-hydrogen has ignored the case.
\vspace{2pt}
\noindent
\begin{tabularx}{\textwidth}{|p{2.2cm}|p{3.0cm}|X|}
  \hline
  \textbf{Written} & \textbf{Maths or} \newline \textbf{science?} &
  \textbf{What it means} \\
  \hline
  $mg$ \rule[-0.65em]{0pt}{2.0em}     & maths & $m$ multiplied by $g$ \\ \hline
  \ch{Mg} \rule[-0.65em]{0pt}{2.0em}  & & \\ \hline
  \ch{CO} \rule[-0.65em]{0pt}{2.0em}  & & \\ \hline
  \ch{Co} \rule[-0.65em]{0pt}{2.0em}  & & \\ \hline
  $xy$ \rule[-0.65em]{0pt}{2.0em}     & & \\ \hline
  \ch{NaOH} \rule[-0.65em]{0pt}{2.0em} & & \\ \hline
\end{tabularx}

% -- E2 ---------------------------------------------------------------------
\qhead[cbgreen]{E2}{The 1 that nobody writes --- in both subjects}

\noindent Write the number that is \textbf{not printed}. The first one is done for you.

\begin{remember}[Same habit, two subjects]
  In \key{maths}, $s$ means $1s$ --- that is why $8s - s = 7s$. In \key{science},
  \ch{H_2O} has nothing after the \ch{O}, so there is \key{one} oxygen atom.
  Both subjects leave the 1 out and expect you to know it is there.
\end{remember}

\vspace{2pt}
\noindent
\begin{tabularx}{\textwidth}{|p{3.2cm}|X|p{2.4cm}|}
  \hline
  \textbf{Written} & \textbf{The question} & \textbf{The number} \\
  \hline
  $8s - s$ \rule[-0.65em]{0pt}{2.0em}   & how many $s$ is the second term? & 1 \\ \hline
  $w + w$ \rule[-0.65em]{0pt}{2.0em}    & how many $w$ is each term? & \\ \hline
  \ch{H_2O} \rule[-0.65em]{0pt}{2.0em}  & how many oxygen atoms? & \\ \hline
  \ch{NaCl} \rule[-0.65em]{0pt}{2.0em}  & how many chlorine atoms? & \\ \hline
\end{tabularx}

% -- E3 ---------------------------------------------------------------------
% The keystone. Counting the atoms in 3CO2 is not LIKE expanding a bracket; it
% is the same operation on the same kind of object, and the grid works on it
% unchanged.
\qhead[cbgreen]{E3}{Counting atoms is expanding brackets}

\noindent Complete the table. The first row is done for you.

\begin{remember}[The same grid, on a formula]
  \ch{3CO_2} means \key{three lots of} \ch{CO_2}. One \ch{CO_2} is one carbon atom
  and two oxygen atoms. So put it in the grid:

  \medskip
  \noindent
  \begin{minipage}[c]{0.42\linewidth}
    \renewcommand{\arraystretch}{1.6}%
    \begin{tabular}{|c|c|c|}
      \hline
      $\times$ & \ch{C}  & $+\;2\,$\ch{O} \\ \hline
      $3$      & $3\,$\ch{C} & $+\;6\,$\ch{O} \\ \hline
    \end{tabular}
  \end{minipage}%
  \begin{minipage}[c]{0.56\linewidth}
    $3(\ch{C} + 2\ch{O}) = \mathbf{3\ch{C} + 6\ch{O}}$

    \smallskip
    Three \ch{CO_2} is \key{3 carbon atoms and 6 oxygen atoms}.
  \end{minipage}
\end{remember}

\vspace{2pt}
\noindent
\begin{tabularx}{\textwidth}{|p{2.2cm}|X|p{2.7cm}|}
  \hline
  \textbf{Formula} & \textbf{Atoms of each element} & \textbf{Atoms in total} \\
  \hline
  \ch{3CO_2} \rule[-0.65em]{0pt}{2.0em}   & 3 carbon, 6 oxygen & 9 \\ \hline
  \ch{2H_2O} \rule[-0.65em]{0pt}{2.0em}   & & \\ \hline
  \ch{4NaCl} \rule[-0.65em]{0pt}{2.0em}   & & \\ \hline
  \ch{2CaCO_3} \rule[-0.65em]{0pt}{2.0em} & & \\ \hline
\end{tabularx}

\vspace{7pt}
\noindent \textbf{(b)} Now write \ch{2H_2O} the way a maths teacher would. Fill the boxes.

\vspace{4pt}
\noindent\hspace*{1.2em}$2\,(\,2\ch{H} \;+\; \ch{O}\,) \;=\; \nbox \ch{H} \;+\; \nbox \ch{O}$

% -- E4 ---------------------------------------------------------------------
\qhead[cbgreen]{E4}{The same mistake, in two subjects}

\noindent Mr Bowen made two mistakes on the same day. In \textbf{maths} he wrote
\; \expr{x + x = x^2}

\vspace{4pt}
\noindent In \textbf{science} he said: \quad
\emph{``\ch{O_2} has two atoms, so it is a compound.''}

\vspace{4pt}
\begin{enumerate}[label=(\alph*), itemsep=6pt]
  \item Write the correct answer to $x + x$. \hfill \blank[4.0cm]
  \item Is \ch{O_2} an element or a compound? \hfill \blank[4.0cm]
  \item These two mistakes are the \textbf{same} mistake. Explain what it is, in
        full sentences.
\end{enumerate}
\qlines{3}

% -- E5 ---------------------------------------------------------------------
\qhead[cbgreen]{E5}{And he made it again}

\noindent A week later Mr Bowen wrote these two lines, one in each book.

\vspace{4pt}
\noindent \hspace*{1.2em}\textbf{Maths:} \quad $4(x + 3) = 4x + 3$

\vspace{3pt}
\noindent \hspace*{1.2em}\textbf{Science:} \quad \ch{6CO_2} is 6 carbon atoms and 2
oxygen atoms.

\vspace{4pt}
\begin{enumerate}[label=(\alph*), itemsep=6pt]
  \item Write the maths line correctly. \hfill \blank[4.6cm]
  \item Write the science line correctly. \hfill \blank[4.6cm]
  \item What did he forget to do, both times? Answer in one sentence.
\end{enumerate}
\qlines{2}

% ===========================================================================
% ANSWER KEY -- TEACHER COPY ONLY
% ===========================================================================
\ifkey
\clearpage

\hwsection[cbred]{Answer Key --- Teacher Copy}

\linespread{1.0}\selectfont
\setlist[enumerate]{itemsep=2pt, topsep=3pt, leftmargin=*}
\setlist[itemize]{itemsep=2pt, topsep=3pt, leftmargin=*}

\qhead[cbred]{A1}{How many terms?}
\noindent $4a+3b$: \textbf{2} \quad $5x$: \textbf{1} \quad $2p+7q+4$: \textbf{3}
\quad $8m-3m+n$: \textbf{3} \quad $6$: \textbf{1} \quad $y+2z-5+3y$: \textbf{4}

\smallskip
\noindent{\small The two rows that separate the class are $6$ and $2p+7q+4$. A number
on its own \emph{is} a term, and students who have decided a term must contain a letter
write 0 for the first and 2 for the second. Both errors come from one belief, so fix it
once.}

\qhead[cbred]{A2}{Like or not like?}
\noindent Left column: $3a$, $5a$ \textbf{like} \sep $4x$, $4y$ \textbf{not} \sep
$7$, $2$ \textbf{like}. \\
Right column: $2ab$, $5ba$ \textbf{like} \sep $m$, $m^2$ \textbf{not} \sep
$5c$, $-2c$ \textbf{like}.

\smallskip
\noindent{\small $2ab$ and $5ba$ is the row worth the time: multiplication can be
written in either order, so they are the same kind. If anyone doubts it, put $a=2$ and
$b=3$ in both. \\
$5c$ and $-2c$ is like. The sign belongs to the term but does not change what kind of
term it is --- which is exactly the fact A4 needs. \\
$m$ and $m^2$ is the only row that cannot be settled with an everyday picture. $m^2$
means $m \times m$, and that is a different thing from $m$.}

\qhead[cbred]{A3}{Simplify}
\noindent (a) $5b$ \quad (b) $5k$ \quad (c) $8m+3n$ \quad (d) $6p+6q$ \quad
(e) $7s$ \\
(f) $7x$ \quad (g) $3a+7$ \textbf{already finished} \quad (h) $7y$ \quad
(i) $3g$ \quad (j) $5c+3d$ \textbf{already finished}

\smallskip
\noindent{\small (e), (f), (h) and (i) all contain a lone letter, and skipping it goes
wrong in \emph{two} directions, which is worth knowing before you mark. Where the lone
letter is \textbf{added}, the answer comes out one short: (f) gives $6x$. Where it is
\textbf{subtracted}, the answer comes out one too many: (h) gives $8y$ and (i) gives
$4g$. (e) is the extreme case --- $8s-s$ becomes $8$, with the letter gone altogether.
Ask for the 1 in pencil until it stops being needed. \\
(g) and (j) are the two that must be left alone, and the instruction line says there are
two. A student who simplifies (g) to $10a$ has not read it; a student who leaves eight
alone has read it too carefully. \\
(h) is the hardest item: the invisible 1 is in the middle of the line \emph{and} carries
a minus sign.}

\qhead[cbred]{A4}{Every sign travels with its own term}
\noindent (a) $9x-5x+4y+2y = \mathbf{4x+6y}$ \quad
(b) $8a+2a-3b-b = \mathbf{10a-4b}$ \\
(c) $5t-2t+6-1 = \mathbf{3t+5}$ \quad
(d) $7p-3p-2q+6q+4 = \mathbf{4p+3q+4}$

\smallskip
\noindent{\small Mark the middle line, not the answer. Every error here is invisible in
a final number and obvious in the working, which is the whole argument for demanding
it. \\
(b) is the one that goes wrong: the $-b$ is a lone letter \emph{and} carries a minus, so
it collects two errors at once and comes out $10a-2b$ or $10a-4$. \\
(d) ends with a number that belongs to nothing, and a class that has just learned to
collect will try to collect it.}

\qhead[cbred]{A5}{Now you write the question}
\noindent No single right answer. Check the four conditions \textbf{in this order}:
\begin{enumerate}[label=\arabic*., itemsep=3pt]
  \item \textbf{Are there two different kinds of thing?} This is the one that fails. A
        question about 5 red pens and 3 blue pens is still one kind of thing being
        counted, and it gives $8p$, not $5a+3b$.
  \item \textbf{Is it said what $a$ and $b$ stand for?}
  \item \textbf{Is each letter a number, not an object?} ``$a$ is a mango'' is the
        error; ``$a$ is the number of mangoes in one bag'' is right.
  \item \textbf{Full sentences, ending in a question?}
\end{enumerate}

\smallskip
\noindent{\small One that works: ``A bag holds $a$ apples. A box holds $b$ bananas.
Mr Bowen buys 5 bags and 3 boxes. Write an expression for the number of pieces of fruit
he has.'' \\
A question ending ``what is $5a+3b$?'' has not done the task --- it is the expression
read out in words. Name that out loud rather than marking it wrong. \\
Read four or five aloud. This is the best item in Section A for finding out who is
reading.}

\qhead[cbred]{B1}{Fill the grid}
\noindent (a) $3m + 12$ \quad (b) $5y + 10$ \quad (c) $8n + 12$

\smallskip
\noindent{\small (c) is the only one where the left-hand box is not a lone letter, and
it is the one to check: $4 \times 2n$ is $8n$, not $6n$ and not $8n^2$. That is B2's
second rule arriving before it is stated, on purpose --- the grid makes it obvious and
the rule can be named afterwards. \\
Strong students will start writing the answer straight out. That is fine; ask them to
draw the grid for (c) anyway.}

\qhead[cbred]{B2}{Expand}
\noindent (a) $3a+6$ \quad (b) $5b+15$ \quad (c) $9y+27$ \quad (d) $8+4f$ \quad
(e) $56+8z$ \quad (f) $6k-18$ \\
(g) $2y-8$ \quad (h) $5m-30$ \quad (i) $12-4t$ \quad (j) $6q+15$ \quad
(k) $20u-4$ \quad (l) $6+12v$ \\
(m) $48+32w-24g$

\smallskip
\noindent{\small (d) and (e) put the number first inside the brackets, and a student who
has learned ``letter then number'' writes $4f+8$. That is the same expression, so accept
it and say why out loud. \\
(i) is the one to watch: the letter is second, the answer starts with a number and
\emph{looks} unfinished --- and B3 is about to say that it is not. \\
(m) has three terms and the third gets left behind. Counting the terms inside before
starting is A1 doing real work.}

\qhead[cbred]{B3}{Knowing when to stop}
\noindent $12-4c$ finished \quad $5x+3x \rightarrow \mathbf{8x}$ \quad
$3a+6$ finished \quad $8+2m+5 \rightarrow \mathbf{2m+13}$

\smallskip
\noindent{\small $8+2m+5$ is the useful row. The two numbers are \emph{not} next to each
other, and a student reading left to right will not see them as a pair. It is A4's skill
in B's clothing. \\
This question is why Section A comes before Section B. A student who can expand
perfectly and cannot do this has learned two lessons and not joined them.}

\qhead[cbred]{B4}{Mr Bowen's homework, again}
\noindent (a) \textbf{wrong} --- $\mathbf{3a+12}$. He multiplied the $a$ and not the 4. \\
(b) \textbf{wrong} --- $\mathbf{8b-20}$. He did $4+5$ instead of $4 \times 5$. \\
(c) \textbf{right}. \\
(d) \textbf{wrong} --- $\mathbf{6m+12}$. He did $6+2$ instead of $6 \times 2$.

\smallskip
\noindent{\small (b) and (d) are the same error and are printed apart on purpose: a
student who corrects one and not the other is guessing. \\
(a) is the headline error of 2.4, and it reappears in E5 wearing a chemical formula. If
you mark one question in Section B, mark this one, and keep the answers to show
alongside E5.}

\qhead[cbred]{B5}{Expand, then collect}
\noindent (a) $6y+3+5y+4 = \mathbf{11y+7}$ \quad
(b) $2a+12+3a+3 = \mathbf{5a+15}$

\smallskip
\noindent{\small The order is the lesson: expand, \emph{then} collect. A student who
tries to collect first meets $2y+1$ inside a bracket and stalls. \\
(b) has two brackets and is the only place in the packet where one expansion has to be
held while another is done. Insist on both being written out before anything is
collected.}

\qhead[cbred]{B6}{The fish tanks}
\noindent (a) $7(f+2)$ \quad (b) $7f+14$ \quad (c) $\mathbf{7f}$

\smallskip
\noindent{\small (c) is the best item in Section B. The answer is $7f$, with no number
added and no bracket --- fourteen rubber ducks simply are not alive. Expect blanks, and
expect $7f+14$. Anyone who writes $7f$ has read every word. \\
An expression that is a single term with no operation in it does not look like an answer
to this class. Say that it is one.}

\qhead[cbred]{C1}{Use the Periodic Table in your book}
\noindent magnesium \ch{Mg}, period 3, group 2, \textbf{metal} \\
chlorine \ch{Cl}, period 3, group 7, \textbf{non-metal}

\smallskip
\noindent{\small The two are in the same period and different groups on purpose: one row
across, two columns apart, so a student who has period and group the wrong way round
gives 2 and 3, and 7 and 3, and is wrong in a way you can read at a glance. \\
This is the only question on the packet that cannot be answered without opening the
book. A blank here may mean the book stayed in the bag rather than that the science is
missing --- ask before marking it. \\
Group numbers follow the book's table, which numbers them 1 to 8. A student using a wall
chart numbered 1 to 18 will say 17 for chlorine. That is a different table, not a wrong
answer; ask which one they used.}

\qhead[cbred]{C2}{Match the word to its meaning}
\noindent 1--D \quad 2--F \quad 3--H \quad 4--G \quad 5--A \quad 6--B \quad 7--E
\quad 8--C

\smallskip
\noindent{\small \emph{Period} (4--G) and \emph{group} (5--A) are the pair that swap, and
not because of the science: both words have an everyday English meaning --- a lesson,
and people together --- and neither helps. Make the class say \textbf{period is a row,
group is a column} out loud. \\
A student with 4 and 5 reversed and everything else right has a vocabulary problem, not
a chemistry problem.}

\qhead[cbred]{C3}{The capital-letter rule}
\noindent (a) \ch{NA} $\rightarrow$ \ch{Na}, sodium \quad
\ch{cl} $\rightarrow$ \ch{Cl}, chlorine \quad
\ch{mG} $\rightarrow$ \ch{Mg}, magnesium \quad
\ch{CA} $\rightarrow$ \ch{Ca}, calcium

\smallskip
\noindent (b) \ch{CO} is \textbf{two} symbols, \ch{C} and \ch{O}, so he has written
carbon and oxygen bonded together --- a compound, not an element. Cobalt is \ch{Co}, with
a small second letter.

\smallskip
\noindent{\small (b) is the question. The capital letter is not neatness, and this is the
one place a student can see that it changes the meaning completely: \ch{CO} and \ch{Co}
are different substances. \\
It is also E1 in advance, so keep the answers.}

\smallskip
\noindent (c) Magnesium is in \textbf{period 3}, \textbf{group 2}, and it is a
\textbf{metal}.

\smallskip
\noindent{\small (c) is the only question on the packet that cannot be answered without
opening the book, so a blank here may mean the book stayed in the bag rather than that
the science is missing. Ask before marking. \\
Group numbers follow the book's table, which numbers them 1 to 8. A student using a wall
chart numbered 1 to 18 will still say 2 for magnesium --- but for carbon or chlorine the
two systems disagree, so if this question is ever widened, say which table to use.}

\qhead[cbred]{C4}{Answer in full sentences}
\noindent \textbf{Yes.} Diamond is made of only one kind of atom --- carbon --- so it is an
element. Diamond and graphite look different because the carbon atoms are \textbf{joined
in different arrangements}, not because they are different elements.

\smallskip
\noindent{\small Discuss this one rather than marking it. The everyday evidence --- one
hard and clear, one soft and grey --- points the wrong way, and a student who says ``no,
they look different'' has reasoned from what they can see. Say that is good science,
then give them the extra fact. \\
Watch for ``an element has one atom''. One \emph{kind} of atom; a gold ring has an
enormous number of them.}

\qhead[cbred]{D1}{Element or compound?}
\noindent Elements: \ch{Na}, \ch{O_2}, \ch{H_2}, \ch{Al}. \quad
Compounds: \ch{NaCl}, \ch{CO_2}, \ch{CaCO_3}, \ch{H_2O}.

\smallskip
\noindent{\small \ch{O_2} and \ch{H_2} are the whole question, and most of the class will
tick compound for both. Two atoms is not two kinds. \\
Do not mark these two rows --- put them on the board and take a hand vote before saying
anything. It is the same misconception as E4, and it is worth meeting twice in one
evening.}

\qhead[cbred]{D2}{Naming a compound}
\noindent (a) sodium chloride \quad (b) magnesium oxide \quad (c) lithium sulfide
\quad (d) calcium carbonate

\smallskip
\noindent (e) \textbf{Calcium sulfide.} Two mistakes: the \textbf{metal must come first}
(calcium, not sulfur), and the \textbf{non-metal takes the -ide ending} (sulfide, not
calcide).

\smallskip
\noindent{\small (d) is the only \emph{-ate}, and the trigger is the oxygen. A student
who writes calcium carbide knows the ending rule and did not notice the oxygen. \\
(e) can be got right for the wrong reason. Ask for both mistakes, named separately.}

\qhead[cbred]{D3}{A compound is nothing like the elements in it}
\noindent When atoms bond, the compound has \textbf{new properties} of its own. Sodium
chloride is not sodium and it is not chlorine --- it is a new substance, and it is safe to
eat. What the elements do on their own tells you nothing about what the compound does.

\smallskip
\noindent{\small This is the one sentence of 2.6 that has to survive the week. Push for
the words \emph{new properties}; ``because it is a compound'' is circular and is what
most of them will write first. \\
A student who brings up chlorine in swimming pools has made a real connection and is
also making the point --- pool chlorine is not the chlorine bonded in salt.}

\qhead[cbred]{D4}{Draw the particles}
\noindent \ch{O_2} two circles of one kind, joined --- \textbf{element}. \quad
\ch{H_2O} one large circle with two small ones joined to it --- \textbf{compound}. \\
\ch{CO_2} one central circle with two of a different kind, one each side ---
\textbf{compound}. \quad
\ch{CH_4} one central circle with four of a different kind around it ---
\textbf{compound}.

\smallskip
\noindent{\small Mark two things and no more: \textbf{are the atoms touching}, and
\textbf{is each kind drawn differently}. Loose separated circles mean a mixture, which is
2.7 and not what any of these are. \\
A \ch{CH_4} with three hydrogens is not a drawing error --- it is the small number not
being read, and it is the same failure E3 tests with a big number in front. Do not mark
the artwork.}

\qhead[cbred]{E1}{$\mathrm{Mg}$ or $mg$?}
\noindent \ch{Mg} --- science --- magnesium, one element \\
\ch{CO} --- science --- carbon and oxygen, two elements \\
\ch{Co} --- science --- cobalt, one element \\
$xy$ --- maths --- $x$ multiplied by $y$ \\
\ch{NaOH} --- science --- sodium, oxygen and hydrogen: \textbf{three} elements

\smallskip
\noindent{\small Mark \ch{CO} and \ch{Co} together. Same two letters, same order, one
capital changed, and the answer goes from two elements to one. Nothing in maths behaves
like that, which is why the question is here. \\
\ch{NaOH} is the hardest row and the one worth the most. \ch{Na} is \emph{one} symbol
made of two letters, so a student who reads four letters as four elements --- or as
nitrogen, oxygen and hydrogen --- has ignored the small \emph{a}. It is the same reading
that \ch{CaCO_3} needs in E3. \\
Expect ``\ch{Mg} means $m$ times $g$'' from at least one student who read the shape and
not the case.}

\qhead[cbred]{E2}{The 1 that nobody writes}
\noindent $w+w$ --- 1 each \quad \ch{H_2O} --- 1 oxygen atom \quad
\ch{NaCl} --- 1 chlorine atom

\smallskip
\noindent{\small This is A3 and the science formulae side by side, and the point is that
they are one habit rather than two. Both subjects leave out the 1; neither warns you. \\
A student who has A3 right and this wrong is not being inconsistent --- they have learned
a maths rule rather than a reading rule. Show them the two rows together.}

\qhead[cbred]{E3}{Counting atoms is expanding brackets}
\noindent \ch{2H_2O} --- 4 hydrogen, 2 oxygen --- 6 atoms \\
\ch{4NaCl} --- 4 sodium, 4 chlorine --- 8 atoms \\
\ch{2CaCO_3} --- 2 calcium, 2 carbon, 6 oxygen --- 10 atoms

\smallskip
\noindent (b) $2(2\ch{H} + \ch{O}) = \mathbf{4}\ch{H} + \mathbf{2}\ch{O}$

\smallskip
\noindent{\small \textbf{This is the question the packet was built around.} The two
numbers in \ch{2H_2O} do different jobs: the small 2 counts hydrogen atoms inside one
molecule, and the big 2 in front multiplies all of it. That is exactly the difference
between the 3 inside $4(x+3)$ and the 4 outside it. \\
\ch{2CaCO_3} is the hard row, because the big 2 has to reach past two symbols to get to
the small 3: $2 \times 3 = 6$ oxygen atoms. A student who writes 3 stopped multiplying
halfway along the bracket, which is B4(a) exactly. \\
(b) is not a trick. Atoms in a formula are counted this way; the brackets are simply not
printed.}

\qhead[cbred]{E4}{The same mistake, in two subjects}
\noindent (a) $x+x = \mathbf{2x}$ \quad (b) \textbf{Element.}

\smallskip
\noindent (c) Both times he counted \textbf{how many} and then answered as though it told
him \textbf{how many kinds}. Two $x$ is still one kind of term, so they collect into
$2x$. Two oxygen atoms are still one kind of atom, so \ch{O_2} is still an element.
\textbf{Two of the same thing is not two different things.}

\smallskip
\noindent{\small \textbf{Discuss this one; do not mark it.} It is the only question in
the packet that asks the class to see two errors from two subjects as one error, and the
wording they reach for will be rough. Accept anything that gets to \emph{same kind},
however it is phrased. ``He thought two meant two different'' has it. \\
If the room is silent, put $x+x$ and \ch{O_2} on the board side by side and ask only:
how many \emph{kinds} are in each? The answer arrives on its own.}

\qhead[cbred]{E5}{And he made it again}
\noindent (a) $4(x+3) = \mathbf{4x+12}$ \\
(b) \ch{6CO_2} is \textbf{6 carbon atoms and 12 oxygen atoms}.

\smallskip
\noindent (c) He multiplied the \textbf{first} thing and stopped. The number in front has
to multiply \textbf{everything} after it --- every term inside the brackets, every atom in
the formula.

\smallskip
\noindent{\small Put B4(a) and this question next to each other when you hand the packet
back. They are the same line written twice, once with a letter in it and once with an
element, and a student who corrected B4(a) confidently and left this one alone has
learned a maths rule rather than a rule. \\
(c) is marked on one word: \textbf{everything}. \emph{Each} and \emph{every} carry the
whole of 2.4, which is why they are in B1's word help.}

\fi

\end{document}
